\section{Conclusion and Plan to the Final Evaluation}
\label{sec:conclusion}

A carefully controlled encoder for CLARITY stops at 0.384 Subtask~2 macro-F1 per model (0.405 as a system), and the evidence places the limit in the model rather than in the labels or the decision rule.
Changing only the backbone to a LoRA-tuned 8B decoder improves both subtasks on all ten paired seeds and lifts the system to 0.543 / 0.746, still below the best multi-call pipeline (0.617) and the human reference (0.684).
The gain is broad rather than concentrated where we predicted, which suggests that a larger pretrained model reads better what kind of response an answer is.

\paragraph{Plan.}
Table~\ref{tab:timeline} gives the remaining work up to the final submission (31 October 2026); none of it has been run.
It tests the remaining explanation, \emph{label meaning}, with a cascade that sends only uncertain items, with the classifier's top candidates, to a larger LLM given label definitions and boundary examples, under the same protocol as above.

\begin{table}[t]
\centering
\small
\begin{tabular}{@{}lp{0.77\columnwidth}@{}}
\toprule
Week & Planned work \\
\midrule
Oct 5--11 & Twelve-epoch Qwen on seeds 3--9 (a ten-seed system); all of train with twelve epochs as a one-change test. \\
Oct 12--18 & Cascade: Qwen's top-$k$ candidates, an open LLM choosing with definitions and boundary examples; prompts designed on the train slice. \\
Oct 19--25 & Sweep the deferred fraction for an accuracy-vs-LLM-calls curve; cost per item. \\
Oct 26--31 & Freeze systems; final ten-seed and two-annotator numbers; report and code release. \\
\bottomrule
\end{tabular}
\caption{Timeline to the final evaluation (all planned).}
\label{tab:timeline}
\end{table}
